\documentclass[aps,pra,reprint,superscriptaddress,nofootinbib,longbibliography,showkeys]{revtex4-2}
\usepackage{amsmath,amssymb,amsthm,mathtools}
\usepackage{booktabs,array}

\newtheorem{theorem}{Theorem}[section]
\newtheorem{lemma}[theorem]{Lemma}
\newtheorem{corollary}[theorem]{Corollary}
\newtheorem{proposition}[theorem]{Proposition}
\newtheorem{conjecture}[theorem]{Conjecture}
\theoremstyle{definition}
\newtheorem{definition}[theorem]{Definition}

\newcommand{\TT}{\mathrm{TT}}
\newcommand{\rank}{\operatorname{rank}}
\newcommand{\im}{\operatorname{im}}
\newcommand{\Mat}{\operatorname{Mat}}
\newcommand{\CDet}{\operatorname{CDet}}
\newcommand{\perm}{\operatorname{perm}}
\newcommand{\FEMPS}{\textnormal{\textsc{femps}}}
\begin{document}
\title{Fermionic Extensions of Continuous Functional Tensor Networks:\\
Exchange-Rank and Exterior-Contraction Obstructions}
\author{Author information withheld for review}
\affiliation{Affiliation withheld for review}
\keywords{fermionic tensor networks, first quantization, exterior algebra,
computational complexity, continuous functional basis}

\begin{abstract}
Hong et al. proposed continuous functional tensor networks by expanding each
particle coordinate in an orthonormal functional basis, representing the
coefficient tensor by a tensor network, evaluating basis-projected operators
by contraction, and optimizing the tensors by automatic differentiation.  We
study the fermionic extension that their work left open.  Two independent
obstructions emerge.  First, an ordinary particle-site tensor train pays an
unavoidable binomial bond for exchange statistics; we recall the known flat
particle spectrum of a single Slater determinant.  Second, moving
antisymmetry into exterior algebra
removes that representation floor but does not make compact inputs easy to
contract: exact squared-norm evaluation is \#P-hard already at fixed one-form
bond two under a polynomial-time metric reduction, and an independent
reduction gives exact point-evaluation hardness in an explicit rational
shifted-Legendre functional basis at the same bond.  Along a matrix-pair
subfamily, bounded and selected fixed-state memories collapse to polynomial
linear combinations of established antisymmetrized geminal powers, while a
bandwidth-one unique path already encodes a permanent.  A generic relative
squared-norm scheme would also approximate real positive-semidefinite
permanents, whereas entrywise-nonnegative and additive promised cases remain
open.  Finally, cut-rank divisibility rejects a universal direct tensor product
of a statistics carrier and a Slater-normalized correlation multiplicity.  A
nonbranching exterior subclass is evaluated as a numerical consequence of
these constraints.  Because that
subclass is algebraically a finite nonorthogonal Slater expansion, the data are
not claimed as a method distinct from nonorthogonal configuration interaction
(NOCI).  These exact obstructions constrain, but do not exclude, genuinely
distinct restricted or controlled approximate algorithms; they do not
establish a scalable generic fermionic solver.
\end{abstract}
\maketitle

\section{Scope and thesis}

The direct motivation is the functional-tensor-network construction of
Hong--Xiao--Hu--Ji--Ran \cite{Hong2022FunctionalTN}.  For orthonormal local
functions $\{\phi_s\}_{s=1}^D$, a first-quantized wavefunction is written
\begin{equation}
 \Psi(x_1,\ldots,x_N)=
 \sum_{s_1,\ldots,s_N}C_{s_1\cdots s_N}
 \prod_{i=1}^N\phi_{s_i}(x_i).
 \label{eq:functional-expansion}
\end{equation}
Orthonormality turns the continuum norm into the Euclidean norm of $C$ and
turns differential and multiplicative operators into finite functional-basis
matrices.  Hong et al. compressed $C$ by an MPS and optimized its Rayleigh
quotient with automatic differentiation; they explicitly named fermionic
tensor networks as a future electronic extension.

For identical fermions, however, $C\in\Lambda^N V_D$.  An ordinary MPS over
the labelled particle slots and an occupation-number MPS use different tensor
product factorizations.  Their bond dimensions are not related by the abstract
Hilbert-space isomorphism.  Hilbert-space MPS \cite{LiChan2016HSMPS},
Grassmann/graded networks \cite{KongZhuXie2026GrassmannTN}, and the ordered
first-quantized construction of Li--Waintal
\cite{LiWaintal2026FirstQuantization} are therefore comparators, not aliases
for the representation studied here.

All complexity statements below use explicit finite encodings and specify the
underlying field.  Floating-point calculations are separated from the formal
results and are reported with their basis, truncation, optimization, and
initialization controls.  Antisymmetry is structural in every exterior or
ordered-chamber construction below, so the corresponding residual is
$\epsilon_{\rm AS}=0$.  The generic exact-contraction obstruction motivates
restricted and controlled approximate algorithms rather than excluding them.

The thesis is deliberately negative and representation-specific:
\begin{quote}
Ordinary particle TT stores exchange multiplicity, whereas compact exterior
cores can hide hard coefficient evaluation.  Removing the former cost does
not remove the latter, and neither observation is a no-go theorem for every
fermionic tensor network, ordered sector, determinant expansion, or promised
approximation.
\end{quote}

\section{Exchange rank in the labelled-particle factorization}

Let $V$ be a finite-dimensional real or complex Hilbert space.  Identify
$\Lambda^N V$ isometrically with the alternating subspace of $V^{\otimes N}$
by
\[
 J_N(v_1\wedge\cdots\wedge v_N)=\frac{1}{\sqrt{N!}}
 \sum_{\pi\in S_N}\operatorname{sgn}(\pi)
 v_{\pi(1)}\otimes\cdots\otimes v_{\pi(N)}.
\]
Write $C_{(k)}$ for the $k\mid N-k$ particle matricization.

\begin{theorem}[Structural result I: exact particle-TT ranks]
For $1\le k<N$, the minimal exact ordinary-TT bond is
\[
 r_k^{\TT}(C)=\rank C_{(k)}.
\]
\label{thm:r1}
\end{theorem}

\begin{proof}
Write an arbitrary TT coefficient at the $k\mid N-k$ cut as
$C_{(k)}=L_kR_k$, where the columns of $L_k$ and rows of $R_k$ are indexed by
the cut bond.  Thus every exact TT has bond at least $\rank C_{(k)}$.  For the
converse, choose a basis of the column space of the first unfolding and express
every remaining slice in that basis.  Repeat on the residual coefficient array
from left to right.  At step $k$ the number of retained basis vectors is
exactly $\rank C_{(k)}$, and after the final step their contracted product is
$C$.  This elementary rank-factorization construction is the exact TT-SVD
argument \cite{Oseledets2011TT}; it proves simultaneous attainability at every
cut and not merely separate lower bounds.
\end{proof}

\begin{theorem}[Structural result II: universal exchange floor]
If $0\ne C\in\Lambda^N V$, then
\[
\rank C_{(k)}\ge \binom{N}{k}.
\]
The lower bound is attained by every nonzero decomposable $N$-form.
\label{thm:exchange-floor}
\end{theorem}

\begin{proof}
Choose a nonzero coefficient $c_I$ on an $N$-element support $I$.  Restrict rows to
$k$-subsets $A\subset I$ and columns to $I\setminus A$.  The resulting
$\binom{N}{k}$ square minor is diagonal up to signs: its diagonal entries are
$\pm c_I$, while an off-diagonal entry repeats an index and vanishes.
A decomposable form supported on an $N$-plane has no larger contraction image.
\end{proof}

\begin{corollary}[strict-antisymmetry truncation floor]
If $\widetilde C\ne0$ and
$r_k^{\TT}(\widetilde C)<\binom{N}{k}$ at any cut, then
$\widetilde C\notin\Lambda^N V$.
\end{corollary}

The TT-rank characterization is standard.  The exterior rank-floor proof is
included for audit, but no priority claim is made pending a dedicated exterior-
algebra literature review.  It is distinct from the stronger, generally open
entropy-minimization statements studied for fermionic reduced density
matrices \cite{CarlenLiebReuvers2016FermionicEntropy,
LiWaintal2026FirstQuantization}.

\begin{theorem}[Structural result III: known flat Slater particle spectrum]
For orthonormal $u_1,\ldots,u_N$, the normalized Slater determinant $\Phi$
has Schmidt decomposition
\[
 \Phi=\frac{1}{\sqrt{\binom{N}{k}}}
 \sum_{|I|=k}(-1)^{\epsilon(I)}\Phi_I\otimes\Phi_{I^c}.
\]
Thus all $\binom{N}{k}$ nonzero singular values equal
$\binom{N}{k}^{-1/2}$.  The best cut-rank-$r$ approximation has relative
squared error
\[
 \epsilon_{k,r}^2=1-\frac{r}{\binom{N}{k}},
\]
so error at most $\epsilon$ requires
$r_k^{\TT}\ge\lceil(1-\epsilon^2)\binom{N}{k}\rceil$.
\label{thm:r2}
\end{theorem}

\begin{proof}
Partition the signed permutation sum by the set of orbitals in the first $k$
slots.  Each class factors into normalized internal Slater sums with coefficient
$\sqrt{k!(N-k)!/N!}$.  Distinct subsets give orthonormal states on both sides.
The approximation formula is the Eckart--Young theorem applied to this flat
singular-value list.  This is the standard particle-reduced-density-matrix
structure of a Slater determinant \cite{Coleman1963FermionRDM}.
\end{proof}

No flat spectrum is asserted for a general correlated alternating tensor.
Theorems~\ref{thm:r1}, \ref{thm:exchange-floor}, and~\ref{thm:r2} concern
labelled-particle TT only; they do
not apply to mode MPS, ordered chambers, or exterior/determinant carriers.

For a fixed orthonormal one-particle basis of size $D$, all of these
representations describe states in exactly the same finite physical space
$\Lambda^N V_D$, whose full-configuration-interaction dimension is
$\binom{D}{N}$.  A mode-space FCI or quantum-chemistry DMRG calculation in
that orbital basis changes the factorization and variational compression, not
the physical finite-basis space \cite{ChanSharma2011DMRG,LiChan2016HSMPS}.
Consequently, \FEMPS{} has neither a larger variational space nor an accuracy
or efficiency advantage merely by being first quantized.  Such an advantage
would require matched basis, tolerance, timing, and memory comparisons.

\section{Exterior matrix-product representation}

\begin{definition}[one-form matrix-wedge state]
For site $j$, let
$A^{[j]}\in\Mat_{\chi_{j-1}\times\chi_j}(V_D)$ with
$\chi_0=\chi_N=1$.  For compatible matrices of forms define
\[
 (X\wedge Y)_{ac}=\sum_b X_{ab}\wedge Y_{bc}.
\]
The open-boundary state is
\begin{equation}
 C(A^{[1]},\ldots,A^{[N]})=
 [A^{[1]}\wedge\cdots\wedge A^{[N]}]_{11}
 \in\Lambda^N V_D.
 \label{eq:femps}
\end{equation}
\end{definition}

Exterior associativity makes matrix-wedge multiplication associative, and each
virtual path is a wedge of $N$ one-forms.  Hence the state is alternating by
construction and $\epsilon_{\rm AS}=0$.  The $\chi=1$ family is exactly the
decomposable forms.  A sum of $R$ Slaters embeds by $R$ diagonal virtual paths,
and the ordinary scalar MPS gauge action remains valid.  These elementary
facts do not supply a canonical form or an efficient norm contraction.  The
displayed parameter count $O(ND\chi^2)$ is an input-size statement only.

\section{Fixed-bond exact contraction obstruction}

For an $n\times n$ matrix $M=(M_{ij})$ over a noncommutative algebra, define
the row-ordered Cayley determinant
\[
 \CDet(M)=\sum_{\sigma\in S_n}\operatorname{sgn}(\sigma)
 M_{1,\sigma(1)}\cdots M_{n,\sigma(n)}.
\]
Chien--Harsha--Sinclair--Srinivasan prove \#P-hardness over
$\Mat_2(\mathbb F)$ in characteristic zero
\cite{ChienHarshaSinclairSrinivasan2011NoncommDet}; related hardness results
for the noncommutative determinant are given by Arvind--Srinivasan
\cite{ArvindSrinivasan2018NoncommDet}.

The distinction from exchange is essential.  For scalar entries ($\chi=1$),
the factors commute, so the same alternating permutation sum is the ordinary
determinant and is computable by Gaussian elimination; in \FEMPS{} it is just
the Slater exchange carrier.  At $\chi=2$, the entries lie in the fixed
four-dimensional algebra $\Mat_2(\mathbb Q)$ and the product must remain in row
order.  The final physical wavefunction is still scalar and exactly
antisymmetric: the noncommutativity lives only in the virtual computation.

\begin{lemma}[direct Cayley coefficient]
Let $M_{ij}\in\Mat_d(\mathbb F)$ and set
$\mathcal A^{[i]}=\sum_jM_{ij}e_j$.  For $u,v\in\mathbb F^d$,
\[
 u^{\mathsf T}\mathcal A^{[1]}\wedge\cdots\wedge
 \mathcal A^{[n]}v
 =\bigl(u^{\mathsf T}\CDet(M)v\bigr)
 e_1\wedge\cdots\wedge e_n.
 \label{eq:direct-cayley}
\]
After absorbing the boundaries into the endpoint cores, this is a FEMPS with
$D=N=n$ and maximum internal bond $d$.
\end{lemma}

\begin{proof}
A nonzero top-form term chooses every physical basis vector once.  The choice
is a permutation, its exterior sign is the permutation sign, and its virtual
matrix product remains in increasing site order.  Their sum is $\CDet(M)$.
\end{proof}

\begin{theorem}[Fixed-bond squared-norm hardness]
Over $\mathbb Q$, exact squared-norm evaluation for rational one-form FEMPS is
\#P-hard under a polynomial-time metric reduction even when $D=N$ and every
internal bond is at most two.  A generic exact contraction algorithm for fixed
$\chi=2$ would imply $\mathrm{FP}=\#\mathrm P$.
\label{thm:c1}
\end{theorem}

\begin{proof}
Let $\varphi$ be a 3CNF instance with $m$ clauses and $S$ satisfying
assignments.  Theorems 3.5 and 3.9 of
Chien--Harsha--Sinclair--Srinivasan (CHSS) construct in polynomial time a
polynomial-size matrix $H_\varphi$ over $\Mat_2(\mathbb Q)$, with constant-bit
integer gadget weights, such that
\begin{align}
 \CDet(H_\varphi)&=aI_2+bJ_2,
 &J_2&=\begin{pmatrix}0&1\\1&0\end{pmatrix},\nonumber\\
 a+b&=4^{3m}S.&&
\end{align}
CHSS label the graph vertices, and their determinant is the row-ordered Cayley
determinant in that vertex order.  Apply Eq.~\eqref{eq:direct-cayley} with
$u=e_1$ and $v=e_1+e_2$.  The sole top-form coefficient is
\[
 u^{\mathsf T}(aI_2+bJ_2)v=a+b=4^{3m}S\ge0,
\]
so one exact squared-norm query returns $16^{3m}S^2$.  Exact nonnegative
integer square root followed by division by $4^{3m}$ recovers $S$.  The CHSS
graph size is polynomial in $|\varphi|$; all FEMPS cores have polynomial
encoding length; and the queried norm and its square root have polynomial bit
length because $S\le2^{n_{\rm var}}$.  Construction and postprocessing are
therefore polynomial-time, and the maximum internal bond remains two.
\end{proof}

The structured CHSS boundary is what sharpens the bond from three to two.  For
an arbitrary signed Cayley entry $x$, one may still adjoin the bond-one
reference $\Theta=e_1\wedge\cdots\wedge e_n$ by a block direct sum.  The two
norms $x^2$ and $(x+1)^2$ recover
$x=((x+1)^2-x^2-1)/2$ at maximum bond three.  That polarization observation is
useful for general outputs, but it is not needed for the CHSS reduction.

CHSS work over any field of characteristic $p\ne2$: their characteristic-zero
statement is \#P-hard and their positive odd-characteristic statement is
$\mathrm{Mod}_p\mathrm P$-hard.  We specialize to $\mathbb Q$ so that the
input, output, square root, and bit-complexity model are explicit; no statement
is made in characteristic two.

The result is not in conflict with the polynomial collapse of a fixed
$\Mat_2$ homogeneous pair power below.  The hard one-form cores are
site-dependent and preserve row order.  Repeating one even two-form instead
symmetrizes factor order.

The direct coefficient identity is pointwise algebraic and also holds when
the entries are one-particle functions evaluated at a coordinate.  A separate
reduction for an explicit rational functional basis is proved in
Theorem~\ref{thm:pointwise}; it is independent of the norm reduction even
though both start from the same CHSS source.  Neither result says that every
$\chi\ge2$ \FEMPS{} is unsuitable for Monte Carlo or controlled approximation.

\subsection{Why parity is only part of the mechanism}

The contrast above is not an odd-versus-even dichotomy.  Site-dependent odd
one-form cores retain the ordered virtual product and can encode the hard
row-ordered determinant.  Repetition of one homogeneous even two-form makes
the physical factors commute and can collapse a fixed coefficient algebra.
But scalar-bond odd one-form states are just Slater determinants and are easy,
whereas the growing-width even sparse-path family below is hard.  Parity
explains the commutation step, while site dependence, homogeneity, virtual
width, and algebraic closure decide the contraction problem.  Splitting a
general matrix-valued pair core into consecutive one-form cores can also
increase the bond by a factor of order $D$; it is not a cost-free equivalence.

\subsection{Independent tagged pair-power mechanism}

For $N=2M$, a matrix-valued two-form
$\Omega_B=\sum_{p<q}B_{pq}e_p\wedge e_q$ defines
\[
 \Psi_M(B;l,r)=l^{\mathsf T}\frac{\Omega_B^M}{M!}r.
\]
Shift tags $E_{i,i+1}\otimes M_{ij}$ annihilate all factor orders except
increasing row order.  This gives a separate polynomial-time Turing reduction
with pair bond $2(n+1)$ and one-form bond $O(N^2)$.  It is weaker than
Theorem~\ref{thm:c1} for unrestricted one-form cores, but isolates growing
noncommutative order memory inside the homogeneous pair-power subfamily.

\section{The matrix-pair corridor}

\subsection{Bounded Wedderburn--radical memory}

\begin{theorem}[Bounded-algebra LC--AGP collapse]
Let $\mathcal A$ be a finite-dimensional unital complex algebra with radical
$R^d=0$ and Wedderburn--Malcev decomposition
\[
 \mathcal A=S\oplus R,\qquad
 S=\bigoplus_{a=1}^q\Mat_{p_a}(\mathbb C).
\]
Put $p=\max p_a$ and $\rho=\dim R$.  For any linear boundary functional
$\lambda$, $\lambda(\Omega^M)/M!$ is an exact LC--AGP with
\begin{align}
 K&\le \sum_{k=0}^{\min(d-1,M)}q^{k+1}\rho^k
 \binom{M+v_k-1}{v_k-1},\nonumber\\
 v_k&=(k+1)p^2+k.
 \label{eq:bounded-collapse}
\end{align}
For fixed $p$ and $d$, this is polynomial in the explicit input dimensions.
If all data, the displayed decomposition, and the boundary are supplied over
$\mathbb Q$ with polynomial-bit structure constants, the LC--AGP coefficients
can be constructed in deterministic polynomial Turing time and have
polynomial bit length.  No algorithm for discovering a rational
Wedderburn--Malcev decomposition is asserted here.
\end{theorem}

\begin{proof}
Expand $\Omega=\Omega_S+\Omega_R$ into noncommutative words.  A word with at
least $d$ radical factors vanishes.  For a surviving number $k$, resolve the
$k+1$ semisimple runs through the $q$ blocks and each radical insertion in a
basis of $R$, giving at most $q^{k+1}\rho^k$ structural choices.  For one
choice, even physical two-forms commute; after applying $\lambda$, the result
is a homogeneous degree-$M$ scalar polynomial in at most
$v_k=(k+1)p^2+k$ two-forms.  In characteristic zero, $M$th powers of linear
forms span $\operatorname{Sym}^M(\mathbb C^{v_k})$.  Each such power divided by
$M!$ is one AGP, proving Eq.~\eqref{eq:bounded-collapse}.  The radical
filtration, block expansion, and polarization count are written out in
Appendix~\ref{app:wm-proof}.
\end{proof}

For full $\Mat_2$, the bound specializes to at most $\binom{M+3}{3}$ AGPs even
though the row-ordered determinant over the same matrix algebra is hard.  This
theorem classifies a pair-power family; it does not give an algorithm for the
site-dependent cores of Theorem~\ref{thm:c1}.

\subsection{Selected growing memories}

\begin{theorem}[Fixed-state graded LC--AGP collapse]
Suppose the coefficient algebra embeds faithfully in
\[
 \Mat_w\!\left(\mathbb C[z_1,\ldots,z_g]/
 (z_1^{d_1},\ldots,z_g^{d_g})\right),
\]
with fixed $w,g$.  Then every boundary contraction of $\Omega^M/M!$ is an
exact LC--AGP with
\[
 K\le\prod_{j=1}^g[M(d_j-1)+1]
 \binom{M+w^2-1}{w^2-1}.
 \label{eq:graded-collapse}
\]
In particular, the one-generator algebra $\mathbb C[z]/z^d$ requires at most
$M(d-1)+1$ AGPs.  For rational inputs with the embedding and quotient basis
supplied explicitly, fixed $w,g$ also gives a deterministic polynomial-time
exact construction with polynomial-bit coefficients.
\end{theorem}

\begin{proof}
Before quotienting, every entry of the $M$th power has degree at most
$M(d_j-1)$ in counter $z_j$.  Tensor-product Vandermonde interpolation
extracts all retained coefficients using the first product in
Eq.~\eqref{eq:graded-collapse}.  At one grid point, a boundary of a $w\times w$
matrix power is a homogeneous degree-$M$ polynomial in $w^2$ commuting
two-forms; polarization supplies the second factor.  Rational nodes give an
exact rational construction for rational inputs; the quotient lift,
interpolation matrices, reduction, and bit bounds are detailed in
Appendix~\ref{app:graded-proof}.
\end{proof}

The theorem includes a noncommutative alternating-word algebra of dimension
$2d-1$ embedded in $\Mat_2(\mathbb C[z]/z^d)$.  Thus growing nilpotency and two
path branches remain insufficient when automaton width and counter count are
fixed.  Weighted-automata and Waring/Veronese ingredients are established
prior art; the project-specific result is only their LC--AGP classification.

\subsection{A bandwidth-one hard path}

Let $P_j=e_{2j-1}\wedge e_{2j}$ be orthonormal pair forms.  Given
$A\in\mathbb F^{M\times M}$, define edge forms
$F_i=\sum_jA_{ij}P_j$ on the unique endpoint path of an upper-bidiagonal
$(M+1)\times(M+1)$ pair matrix.

\begin{theorem}[Sparse APG permanent obstruction]
Over $\mathbb Q$ (or $\mathbb R$ with its positive inner product), for a
zero--one matrix $A$ the endpoint state and its squared norm are
\begin{align}
 \Psi_A&=\frac{F_1\wedge\cdots\wedge F_M}{M!}
 =\frac{\perm(A)}{M!}P_1\wedge\cdots\wedge P_M,
 \nonumber\\
 \langle\Psi_A,\Psi_A\rangle&=\frac{\perm(A)^2}{(M!)^2}.
\end{align}
Exact squared-norm evaluation is therefore \#P-hard for a family with pair
bandwidth one, one virtual path, $D=2M$, and $O(M^2)$ binary data.  Over
$\mathbb C$, the norm formula is $|\perm(A)|^2/(M!)^2$ and the same zero--one
reduction applies.
\end{theorem}

\begin{proof}
Repeated $P_j$ factors vanish.  Every surviving multilinear term selects a
permutation, and even pair forms commute, so its coefficient is the permanent.
For a zero--one matrix, multiply the exact norm by $(M!)^2$ and take the
nonnegative integer square root; Valiant's permanent theorem applies
\cite{Valiant1979Permanent}.
\end{proof}

The state is an established antisymmetrized product of geminals, and permanent-
valued APG/APIG coefficients and APG-to-AGP decompositions have direct prior
art \cite{KawasakiNakatani2024LowRankAPG,
RicherKimAyers2025GraphicalGeminals}.  The hard output is projectively a single
top-form Slater, so this is an input-to-coefficient obstruction, not a lower
bound on physical LC--AGP rank.

\section{Relative approximation and energy certification}

Exact permanent hardness alone does not exclude approximation.  In particular,
entrywise-nonnegative permanents admit an FPRAS
\cite{JerrumSinclairVigoda2004PermanentFPRAS}.  The same sparse-path identity also
accepts a different promise.

\begin{theorem}[Real-PSD relative-norm transfer]
Unless $\mathrm{RP}=\mathrm{NP}$, there is no randomized polynomial-time
relative approximation scheme for the squared norm of every rational real
upper-bidiagonal matrix-pair input, polynomial jointly in dimensions, input
bit length, inverse tolerance, and logarithmic inverse failure probability.
\end{theorem}

\begin{proof}
For rational real positive-semidefinite $A$, the sparse-path construction gives
\[
s_A=\frac{\perm(A)^2}{(M!)^2},\qquad \perm(A)\ge0.
\]
A relative estimate
$\widetilde s_A$ would make $M!\sqrt{\widetilde s_A}$ a relative permanent
estimate after a polynomial tolerance conversion.  Meiburg excludes a PRAS
for real-PSD permanents unless $\mathrm{RP}=\mathrm{NP}$
\cite{Meiburg2023PSDPermanent}.
\end{proof}

Additive permanent estimators remain possible
\cite{AaronsonHance2014Gurvits}.  They do not uniformly imply relative norm
control: cancellations can make the true coefficient exponentially small in
input precision while the additive scale remains order one.

\begin{lemma}[Rayleigh-quotient certificate]
If $n>0$, $E=h/n$, and estimates obey
\[
 |n-\widetilde n|\le\Delta_n,\qquad
 |h-\widetilde h|\le\Delta_h,\qquad
 \widetilde n>\Delta_n,
\]
then, with $\widetilde E=\widetilde h/\widetilde n$,
\[
 |E-\widetilde E|\le
 \frac{\Delta_h+|\widetilde E|\Delta_n}
 {\widetilde n-\Delta_n}.
\]
For random bounds this holds on the simultaneous success event.
\end{lemma}

\begin{proof}
Use
$E-\widetilde E=[(h-\widetilde h)-\widetilde E(n-\widetilde n)]/n$
and $n\ge\widetilde n-\Delta_n>0$.
\end{proof}

Thus estimator unbiasedness alone is not an energy certificate; the
denominator interval must stay positive.  Tensor-network Monte Carlo
\cite{Ferris2015TNMC} and sign-problem complexity
\cite{TroyerWiese2005SignProblem} provide context but are not used as the
reduction.

\section{Failure of a universal direct statistics carrier}

For $0\ne C\in\Lambda^N V$, define the cut contraction
\[
 c_{C,k}:\Lambda^kV^*\longrightarrow\Lambda^{N-k}V,
 \qquad c_{C,k}(\alpha)=\iota_\alpha C,
\]
and $\mathcal B_k(C)=\im c_{C,k}$.  The tested proposal was an exact universal
factorization
\begin{equation}
 \mathcal B_k(C)\cong
 \mathcal S^{\rm fermion}_{N,k}\otimes\mathcal M_k(C),
 \qquad \dim\mathcal M_k(\text{Slater})=1.
 \label{eq:carrier-product}
\end{equation}

\begin{theorem}[Cut-rank divisibility obstruction]
For every $N\ge3$, Eq.~\eqref{eq:carrier-product} cannot hold for every
nonzero $N$-form.
\end{theorem}

\begin{proof}
At the $1\mid N-1$ cut, a Slater determinant has rank $N$, forcing
$\dim\mathcal S^{\rm fermion}_{N,1}=N$.  Hence every admitted rank would be
divisible by $N$.  In $V=\mathbb F^{N+2}$, let
\[
 C_N=e_1\wedge(e_2\wedge e_3+e_4\wedge e_5)
 \wedge e_6\wedge\cdots\wedge e_{N+2}.
\]
Its $N+2$ one-index contractions have distinct monomial supports and are
linearly independent, so $\rank c_{C_N,1}=N+2$.  But $N\nmid N+2$.
\end{proof}

The counterexample is only two Slaters.  Shared physical covectors lock
$N-2$ channels that an independent multiplicity tensor would treat as free.
The rank is invariant under orbital basis changes and direct embeddings, and
full one-cut rank is stable on an open neighborhood.

Standard symmetry-adapted tensor networks remain valid
\cite{SinghPfeiferVidal2010SymmetryTN,Weichselbaum2012NonAbelianTN}.
Particle permutation supplies the one-dimensional sign representation, not a
$\binom{N}{k}$ degeneracy.  Under $GL(V)$, $\Lambda^N V$ is irreducible and the
exterior cut has Littlewood--Richardson multiplicity one while retaining the
full orbital representation.  Hamiltonian-specific charges and multiplets are
not ruled out.  State-adaptive Slater/secant descriptions are established and
need not be canonical \cite{SchliemannCiracKusEtAl2001SlaterRank,
BallicoBernardiCatalisanoChiantini2013GrassmannSecants,
GalganoStaffolani2024GrassmannianIdentifiability}.

The sparse-path construction sharpens the algorithmic warning: whenever its output is nonzero it is
projectively one Slater, yet its scale is a permanent.  A state-intrinsic
multiplicity-one label cannot by itself compile compact cores into a norm
algorithm.

\section{Established positive controls and prior-art boundary}

Fixed-number AGP is equivalent to number-projected pairing structure; projected
BCS, HFB overlaps, and Pfaffian formulas are established
\cite{DukelskyPittelEsebbag2016PBCS,Robledo2009HFBOverlap,
BertschRobledo2012SymmetryRestoration}.  AGP norms and
reduced density matrices have polynomial algorithms
\cite{KhamoshiHendersonScuseria2019AGPRDM}.  Deterministic linear combinations
of optimized and nonorthogonal AGPs precede this project
\cite{UemuraKasamatsuSugino2015AGPCI,
DuttaChenHendersonScuseria2021NonorthogonalAGP}; Pfaffian variational Monte
Carlo with symmetry projection is also established
\cite{TaharaImada2008mVMC,BajdichMitasWagnerSchmidt2008PfaffianQMC}.

Accordingly, finite LC--AGP is used here as a validated tractable control, not
as a new ansatz.  The two collapse theorems above classify when matrix-pair
memory reduces to that control.  Grassmann signs or parity bookkeeping enforce
fermionic algebra but do not, without additional closure, refute either
hardness construction.

\begin{lemma}[Explicit AGP embedding in one-form \FEMPS{}]
Let a scalar two-form be in skew-normal form
$\Omega_F=\sum_{a=1}^r w_a u_a\wedge v_a$.  For $1\le M\le r$,
\[
 \frac{\Omega_F^M}{M!}
 =\sum_{a_1<\cdots<a_M}\left(\prod_{j=1}^M w_{a_j}\right)
 u_{a_1}\wedge v_{a_1}\wedge\cdots\wedge
 u_{a_M}\wedge v_{a_M}
\]
has an exact one-form \FEMPS{} with every internal bond at most
$r\le\lfloor D/2\rfloor$.
\end{lemma}

\begin{proof}
Use the channel label $a$ as the virtual state.  The first core emits
$w_a u_a$, the next core emits $v_a$ diagonally, and every later $u$ core has
only strictly upper-triangular transitions $a<b$ with entry $w_bu_b$; the
following $v$ core is diagonal.  Contracting the final $v$ core with the right
boundary leaves exactly the strictly increasing channel sequences displayed
above.  Thus the construction has $\binom rM$ nonzero paths but bond $r$, and
never enumerates the paths in its core representation.
\end{proof}

For odd particle number, a blocked orbital $b$ gives
$b\wedge\Omega_F^M/M!$ by prepending one bond-one core; no auxiliary
subtraction is needed for this representation statement.  This embedding is
a prior-art control, not a new AGP ansatz.  It also does not contradict the
$O(\chi D)$ bond that can arise when a general matrix-valued pair core of bond
$\chi$ is split into one-form sites: skew-normal scalar AGP structure is the
special closure used by the lemma.

\section{Numerical validation and algorithm-design consequences}

The structural results were subsequently used to choose, rather than reject,
a restricted numerical control.  A nonbranching diagonal-path subclass is
exactly a $K$-term nonorthogonal Slater expansion and evaluates determinant
transition quantities in $O(K^2)$ path pairs without enumerating virtual
histories or materializing the alternating coefficient tensor.  This places it
within established NOCI methodology rather than defining a new correlated
ansatz \cite{ThomHeadGordon2009NOCI,BeylkinMohlenkampPerez2008SlaterSums}.

We retain one representative interacting result.  For six one-dimensional
fermions with harmonic confinement and softened Coulomb repulsion, the
$D=12$, $K=4$ calculation occupies the $\binom{12}{6}=924$ dimensional exterior
space and gives energy $25.0494711446$, compared with same-basis CI
$25.0493664161$.  The error is $1.0473\times10^{-4}$, the energy variance is
$1.1223\times10^{-3}$, the norm error is $1.55\times10^{-15}$, and the
structural antisymmetry residual is exactly zero.  The run took $107.7$ s on
CPU and its sampled peak resident memory was $0.975$ GB.  The $D=12$, $K=4$
orbitals were \emph{not} initialized from a clean $D=12$ Slater: they were
obtained by exact zero-padding of a preoptimized $D=10$, $K=4$ state before the
$D=12$ optimization.  Direct CI was used only as an independent reference.

These floating-point data are a bounded NOCI-equivalent numerical exercise.
They do not establish an ansatz beyond NOCI, an explicit-correlation advantage,
improved functional-basis $D$ convergence, asymptotic scalability, or
superiority over FCI, particle TT, Li--Waintal, or same-basis DMRG.

\section{Discussion}

\subsection{Scope of the complexity statements}

The exact reductions and approximation statements answer different questions.
Exact \#P-hardness does not rule out an FPRAS, an additive estimator, or an
efficient algorithm under a promise that excludes the hard instances.  Norm
polarization recovers an affine signed scale, whereas the sparse-path output is
projectively one Slater; a simple normalized state need not be easy to
normalize from compact cores.  All input-complexity statements count the
explicit rational core data and coefficient bit length, not only the intrinsic
Slater rank of the represented state.

The representation boundary is equally important.  The exchange-rank results
concern labelled-particle TT, while the fixed-bond hardness result concerns
exterior one-form cores.  Neither transfers automatically to mode MPS,
ordered sectors, or every fermionic tensor network.  Likewise, the
matrix-pair collapse repeats one homogeneous even form, whereas the Cayley
construction uses site-labelled one-form cores whose products retain row
order.  Finally, controlled energy estimation requires simultaneous numerator
and norm bounds with a strictly positive lower norm bound; separate unbiased
estimators do not by themselves certify their ratio.

\subsection{Two mechanisms in structured antisymmetrization}
\label{sec:two-mechanism-map}

Consider exact evaluation at a fixed configuration $x=(x_1,\ldots,x_N)$ of
\begin{equation}
 (\mathcal AF)(x)=\sum_{\pi\in S_N}\operatorname{sgn}(\pi)
 F(x_{\pi(1)},\ldots,x_{\pi(N)}).
 \label{eq:structured-antisymmetrizer}
\end{equation}
Two independent properties organize several familiar constructions.  Property
P1 (\emph{assignment factorization}) asks whether, after fixing $x$, every
permutation term has the form
\begin{equation}
 F_\pi(x)=\prod_{i=1}^N a_{i,\pi(i)}(x),
 \label{eq:p1}
\end{equation}
where the permutation selects factors only through $(i,\pi(i))$.  The factors
$a_{ij}(x)$ may depend on the \emph{entire} configuration, as in backflow or
FermiNet; P1 is not a one-particle locality assumption.  Property P2 asks
whether these factors multiply in one commutative coefficient algebra.  In
every cell below we additionally assume that each factor is itself evaluable
from its encoded input in polynomial time and polynomial bit complexity.
Without this base assumption, an easy determinant layer need not make the
complete model easy.

\begin{table*}[t]
\centering
\small
\renewcommand{\arraystretch}{1.12}
\begin{tabular}{@{}l l l@{}}
\toprule
 & P2: commutative & not P2: noncommutative \\
\midrule
P1 & \parbox[t]{72mm}{Slater, generalized Slater, backflow, and FermiNet
determinant channels: the $N!$ Leibniz terms collapse to ordinary
determinants.} & \parbox[t]{72mm}{One-form \FEMPS{}: exterior signs accompany
row-ordered virtual products, giving a Cayley determinant.} \\
not P1 & \parbox[t]{72mm}{The signed Hamilton-path polynomial in
Eq.~\eqref{eq:signed-path}: open cell in the present bounded search.} &
\parbox[t]{72mm}{General explicit antisymmetrization; conditional hardness
only for a precisely defined class that exactly contains the hard bond-two
one-form cores with polynomial overhead.} \\
\bottomrule
\end{tabular}
\caption{A two-mechanism map, not a complete dichotomy or classification
theorem.  P1 concerns assignment structure; P2 concerns the coefficient
algebra.}
\label{tab:two-mechanism-map}
\end{table*}

In the upper-left cell, the ordinary Slater determinant replaces an explicit
$N!$-term Leibniz sum by one determinant.  At fixed $x$, generalized Slater,
backflow, and FermiNet orbital entries remain scalars in a commutative field,
even when every entry depends on all electronic coordinates; permutation
equivariance makes particle exchange act by row or column exchange.  Each
determinant channel is therefore polynomial-time once its entries are
available.  Gaussian elimination applies over fields.  Over a general
commutative ring, divisions need not exist, but division-free determinant
algorithms such as Berkowitz's still give polynomial-size arithmetic circuits
\cite{Berkowitz1984Determinant}.  Hutter's generalized-Slater universality
analysis, FermiNet, the explicitly antisymmetrized layers of
Lin--Goldshlager--Lin, and the separation of Zweig--Bruna concern
representation, universality, ansatz design, or expressivity
\cite{Hutter2020Antisymmetric,Pfau2020FermiNet,
LinGoldshlagerLin2023Antisymmetrized,ZweigBruna2022Separation}; they are not
attributed the complexity map in Table~\ref{tab:two-mechanism-map}.

The upper-right cell isolates the second mechanism.  In one-form \FEMPS{},
Eq.~\eqref{eq:p1} holds with matrix-valued factors, but row order prevents their
exchange.  The exterior algebra supplies $\operatorname{sgn}(\pi)$ while the
virtual algebra retains the noncommuting product, producing the row-ordered
Cayley determinant rather than a scalar determinant.  The following
point-evaluation theorem makes this distinction explicit without changing the
squared-norm theorem above.

\begin{theorem}[Exact rational-polynomial point evaluation]
\label{thm:pointwise}
Let $\ell_r(t)=P_r(2t-1)$, $r\ge0$, be the unnormalized shifted Legendre
polynomials, where $P_r(1)=1$.  Exact pointwise evaluation of rational
one-form \FEMPS{} represented in the basis
$\ell_0,\ldots,\ell_{D-1}$ is \#P-hard under polynomial-time metric
reductions, already when $D=N$ and every internal bond is at most two.
For the isometric convention $J_N$, the exact output is encoded as the
rational $q$ in $\Psi(x)=q/\sqrt{N!}$; equivalently one may request the
unnormalized rational alternating value $q$.
\end{theorem}

\begin{proof}
The shifted polynomials have the integer expansion
\begin{equation}
 \ell_r(t)=\sum_{k=0}^r(-1)^{r-k}
 \binom rk\binom{r+k}{k}t^k
 \label{eq:shifted-legendre}
\end{equation}
and leading coefficient $\binom{2r}{r}$.  Put $t_i=i/(N+1)$ and
$B_{ij}=\ell_{j-1}(t_i)$.  Degree grading and the Vandermonde determinant give
\begin{equation}
 \det B=\left(\prod_{r=0}^{N-1}\binom{2r}{r}\right)
 \prod_{1\le i<k\le N}(t_k-t_i)\ne0.
 \label{eq:legendre-determinant}
\end{equation}
The nodes use $O(\log N)$ bits.  Equation~\eqref{eq:shifted-legendre} shows
that all evaluations use polynomially many bits.  More explicitly,
$Q=(N+1)^{N-1}$ clears every denominator in $B$; the entries of $QB$ have
polynomial bit length, and Hadamard's bound gives polynomial bit length for
its determinant and all minors.  Thus $\det B$, $(\det B)^{-1}$, and
$B^{-1}=Q\,\operatorname{adj}(QB)/\det(QB)$ have polynomial bit length and
are computable by exact rational arithmetic in polynomial time.

For a 3CNF formula $\varphi$, take the CHSS matrix
$H_\varphi=(H_{ij})\in\Mat_N(\Mat_2(\mathbb Q))$ from the proof of
Theorem~\ref{thm:c1}, and define the functional cores
\begin{equation}
 \mathcal A^{[i]}(t)=\sum_{j=1}^N H_{ij}\ell_{j-1}(t),
 \qquad u=e_1,\quad v=e_1+e_2.
 \label{eq:legendre-cores}
\end{equation}
This mapping merely reuses the rational CHSS entries as functional-basis
coefficients, so its size and coefficient bit length are polynomial.  Expanding
the alternating point value at $(t_1,\ldots,t_N)$, and using commutativity only
for the scalar $B_{ij}$, gives
\begin{align}
 \Psi(t_1,\ldots,t_N)
 &=\frac{\det B}{\sqrt{N!}}
 u^{\mathsf T}\CDet(H_\varphi)v \\
 &=\frac{\det B}{\sqrt{N!}},4^{3m}\#\mathrm{SAT}(\varphi).
 \label{eq:legendre-point-value}
\end{align}
Indeed, a repeated basis-column choice vanishes, while a column permutation
$\sigma$ contributes $\operatorname{sgn}(\sigma)\det B$ and leaves the
row-ordered product $H_{1\sigma(1)}\cdots H_{N\sigma(N)}$.  Thus the known
nonzero factor is $c_N=1/\sqrt{N!}$ for $J_N$ and $c_N=1$ for the unnormalized
antisymmetrizer.  The normalized output convention stores the exact algebraic
number as $(q,N)$ meaning $q/\sqrt{N!}$; $N!$ has $O(N\log N)$ bits, so this is
an explicit polynomial-size output domain rather than a hidden algebraic-number
oracle.  One query, followed by rational division by $4^{3m}\det B$, recovers
$\#\mathrm{SAT}(\varphi)$.  This is a polynomial-time metric reduction.
\end{proof}

Theorem~\ref{thm:pointwise} and Theorem~\ref{thm:c1} are independent reductions
from the same Cayley source; neither is inferred from the other.  They concern
worst-case exact evaluation and do not preclude efficient approximate QMC/VMC,
promised subclasses, or particular low-bond instances.

For the lower-left cell, the exact path object is
\begin{equation}
 H_{\rm sgn}(W)=\sum_{\pi\in S_N}\operatorname{sgn}(\pi)
 \prod_{i=1}^{N-1}w_{\pi(i),\pi(i+1)}.
 \label{eq:signed-path}
\end{equation}
Its scalar factors commute, but adjacent vertices couple successive
permutation positions, so it need not have P1.  This is a path, not a closed
cycle: summing a cycle over all permutations would cancel for even $N$ because
a cyclic shift reverses permutation parity.  Hamilton path/cycle polynomials,
sign imbalance, even-minus-odd path counts, fermionants, and immanants were
checked in a bounded search; the latter neighboring families have their own
complexity results \cite{MertensMoore2013Fermionant}.  In the present bounded
literature search, the authors did not find an existing complexity
classification for this exact weighted polynomial.  We therefore register
the cell as open and make no hardness or tractability claim.

The lower-right entry supports only a conditional inheritance statement.  If
a precisely defined parameterized function class exactly represents every
hard bond-two instance in Theorem~\ref{thm:pointwise} with polynomial overhead,
then its exact point-evaluation problem inherits that \#P-hardness.  Universal
approximation does not imply exact containment.  In particular,
Lin--Goldshlager--Lin motivate generic explicit antisymmetrization, including
its factorial sum, but no exact containment lemma is asserted here for their
neural class.

Finally, if $J$ is fully permutation-symmetric, then directly from
Eq.~\eqref{eq:structured-antisymmetrizer},
\begin{equation}
 J\,\mathcal A(M)=\mathcal A(JM).
 \label{eq:jastrow-commutes}
\end{equation}
If $J$ has a compatible scalar particle-TT representation of bond $\chi_J$
and multiplication closes in the chosen functional basis or in a stated
finite extension, tensoring its virtual indices with a Slater carrier gives a
conditional Slater--Jastrow embedding with the corresponding $\chi_J$ bond
factor.  This is an embedding statement, not a norm-hardness reduction; the
point value remains $J(x)\det\Phi(x)$ up to normalization.  The
Zweig--Bruna separation concerns expressive efficiency relative to finite
Slater sums, not point- or norm-evaluation complexity.

The table is an organizing framework and partial classification program.  A
complete complexity classification of structured antisymmetrization,
including the unresolved cells and internal tractability boundaries, is left
for future work.

\subsection{Limitations and next questions}

The results do not establish that every exterior family is hard, that every
tractable exterior family is LC--AGP, or that additive approximation with a
certified polynomial norm lower bound is impossible.  Promised positive,
low-rank, strongly orthogonal, Pfaffian/Gaussian, integrable, or low-treewidth
families remain admissible if their promise excludes the hard embeddings and
their state and energy errors are certified.  A genuinely different
categorical carrier is not excluded by the divisibility obstruction; it must
state its exactness,
reconstruction cost, and contraction theorem rather than assume a free direct
product.

The structural results and all currently admitted restricted-solver numerics
belong to this single manuscript.  No separate \FEMPS{} method paper is claimed
from the diagonal-path data.  A future method manuscript is opened only after
a representation beyond finite NOCI demonstrates a matched, reproducible
advantage---for example, explicit correlation that improves functional-basis
$D$ convergence, or a controlled comparison with Li--Waintal and same-basis
DMRG that identifies a genuine accuracy, stability, or complexity tradeoff.
Higher-dimensional alternating-form rank classification is not pursued unless
a specific rank question determines such an algorithmic or physical decision.

\section{Conclusion}

The attempted fermionic completion of the 2201 functional tensor network has
identified two costs that must not be conflated.  Ordinary particle TT pays a
binomial exchange bond even for a Slater determinant.  Exteriorizing the state
removes that representation charge, but generic exact squared-norm contraction
and rational-polynomial point evaluation are already hard at fixed one-form
bond, while a direct universal carrier
factorization fails dimensionally.  The tested tractable matrix-pair memories
reduce to established LC--AGP structure, and a unique sparse path already
contains a permanent.  The nonbranching $K$-path calculation is correspondingly
a useful NOCI-equivalent numerical control, not yet an independent method
contribution.  The resulting map is an algorithm-design constraint, not a
denial of numerical \FEMPS{}: genuinely distinct restricted subclasses and
error-controlled approximate methods remain open and must be judged by matched
interacting benchmarks, $D$ and correlation convergence, cost, uncertainty,
and antisymmetry residuals rather than by parameter counts alone.

\section*{Use of artificial-intelligence tools}

OpenAI Codex (GPT-5 family, accessed in August--September 2026) assisted with
literature discovery, code and proof exploration, numerical cross-checks, and
language editing under human direction.  The human author selected the
research questions, inspected the cited sources, verified the mathematical
and numerical results, revised the manuscript, and accepts responsibility for
its content.  The AI system is not an author.

\appendix
\section{Detailed bounded-algebra collapse proof}
\label{app:wm-proof}

This appendix expands the algebraic steps used in the bounded-algebra collapse
theorem.  Let $E=\Lambda^{\mathrm{even}}V_D$.  The wedge product makes $E$ a
commutative $\mathbb C$-algebra, although it has zero divisors.  Write
$\Omega=\Omega_S+\Omega_R\in\mathcal A\otimes E_2$, where $E_2=\Lambda^2V_D$.

\paragraph{Radical filtration.}
Because $R$ is a two-sided ideal, a product containing $k$ factors from $R$,
with arbitrary elements of $S$ between them, belongs to $R^k$.  This follows
inductively: $SR^j,R^jS\subseteq R^j$, and
$R^iR^j\subseteq R^{i+j}$.  Hence every word in the expansion of $\Omega^M$
with $k\ge d$ radical factors vanishes.  For fixed $k<d$, every surviving word
has a unique run-length description
\[
 \Omega_S^{m_0}\Omega_R\Omega_S^{m_1}\cdots
 \Omega_R\Omega_S^{m_k},\qquad
 m_0+\cdots+m_k=M-k,
\]
where zero-length semisimple runs are allowed.  Summing over all nonnegative
run lengths is equivalent to retaining the homogeneous degree-$M$ part of the
commutative physical-form polynomial constructed below; no additional
combinatorial factor is omitted because the virtual product order remains the
displayed order.

\paragraph{Wedderburn blocks and radical basis.}
Let $1_a$ be the central identity of the block
$S_a\cong\Mat_{p_a}(\mathbb C)$ and choose a basis
$r_1,\ldots,r_\rho$ of $R$.  Insert
$1_S=\sum_{a=1}^q1_a$ independently into each of the $k+1$ semisimple runs and
expand each radical factor as
\[
 \Omega_R=\sum_{b=1}^{\rho}r_b\,\eta_b,
 \qquad \eta_b\in E_2.
\]
Thus the contribution with $k$ radical occurrences is a sum over at most
$q^{k+1}\rho^k$ choices
$\tau=(a_0,b_1,a_1,\ldots,b_k,a_k)$.  This is an upper bound: incompatible
block transitions may vanish and repeated basis choices may combine.

Fix one $\tau$.  In matrix units of $S_{a_j}$, the corresponding projected
semisimple form has at most $p_{a_j}^2\le p^2$ scalar two-form entries
$x_{j,\mu\nu}\in E_2$.  Together with the selected radical forms
$\eta_{b_1},\ldots,\eta_{b_k}$, the boundary value under $\lambda$ is a
homogeneous degree-$M$ polynomial
\[
 P_\tau(x_{0,\mu\nu},\ldots,x_{k,\mu\nu},
        \eta_{b_1},\ldots,\eta_{b_k})\in E_{2M}
\]
in at most
$v_k=\sum_{j=0}^{k}p_{a_j}^2+k\le(k+1)p^2+k$ commuting two-forms.  Matrix
multiplication and $\lambda$ affect only the scalar coefficients of this
polynomial.  Nilpotence inside $E$ can identify or kill monomials, but cannot
increase the number of powers needed below.

\paragraph{Polarization into AGPs.}
Let $z_1,\ldots,z_v$ be formal commuting variables.  Over characteristic zero,
put $T_M=\{t\in\mathbb Z_{\ge0}^{v-1}:|t|\le M\}$.  The integer simplex is
unisolvent for polynomials of total degree at most $M$: ordered by total degree,
the falling-factorial polynomials
$\prod_{i=2}^v\binom{x_i}{\alpha_i}$ give a triangular evaluation matrix on
$T_M$ with diagonal one.  Hence the square matrix relating the
$|T_M|=\binom{M+v-1}{v-1}$ powers
\[
 (z_1+t_2z_2+\cdots+t_vz_v)^M,\qquad t\in T_M,
\]
to the degree-$M$ monomial basis is nonsingular (the multinomial factors are
nonzero in characteristic zero).  Therefore every homogeneous degree-$M$
polynomial in $v$ variables is a linear combination of at most $|T_M|$ such
powers.

For rational data this is also a Turing construction.  The evaluation entries
are integers of bit length $O(M\log(M+1))$; for fixed $v$ the matrix order is
$O(M^{v-1})$.  Hadamard's inequality and Cramer's rule bound every numerator
and denominator in its inverse by a polynomial number of bits.  Exact Gaussian
elimination therefore computes the power coefficients in polynomial time.  In
the theorem, $v\le v_k$ is fixed when $p,d$ are fixed; the supplied rational
block decomposition and radical basis have polynomial-bit entries, and the
number of structural choices is polynomial.  Multiplication, application of
$\lambda$, and division by $M!$ increase bit lengths only polynomially.

Substitute the scalar two-forms for the $z_i$ and apply the algebra homomorphism
from the symmetric algebra to $E$.  Every resulting term $L^M/M!$ is an AGP.
For a fixed $k$, multiplying the power count by the at most
$q^{k+1}\rho^k$ structural choices and summing over
$0\le k\le\min(d-1,M)$ gives Eq.~\eqref{eq:bounded-collapse}.  All steps use
only the fixed Wedderburn decomposition, radical basis, matrix multiplication,
and exact linear algebra.  For fixed $p$ and $d$, the displayed count is
polynomial in $M,q,$ and $\rho$, hence in the explicit algebraic input size.

\section{Detailed fixed-state graded collapse proof}
\label{app:graded-proof}

Let
$R_{\boldsymbol d}=\mathbb Q[z_1,\ldots,z_g]/
(z_1^{d_1},\ldots,z_g^{d_g})$ and suppose the embedding into
$\Mat_w(R_{\boldsymbol d})$ is part of the input.  Represent every residue by
its canonical polynomial lift with $0\le\deg_{z_j}<d_j$.  An entry of the
unreduced product $\Omega^M$ then has
$\deg_{z_j}\le L_j=M(d_j-1)$.  For each counter use the rational grid
$0,1,\ldots,L_j$.  The univariate Vandermonde matrix is nonsingular, and the
tensor product of these matrices recovers every coefficient of the lifted
product.  Reducing modulo $z_j^{d_j}$ simply discards coefficients whose
$z_j$ exponent is at least $d_j$; it cannot create additional terms.

At one grid point, $\Omega$ is a $w\times w$ matrix whose $w^2$ entries are
commuting physical two-forms.  After the fixed boundary is applied,
$\Omega^M$ is a homogeneous degree-$M$ polynomial in at most $w^2$ such forms.
The integer-simplex polarization construction in
Appendix~\ref{app:wm-proof} expresses it as at most
$\binom{M+w^2-1}{w^2-1}$ powers.  Recombining the at most
$\prod_j(L_j+1)$ grid contributions gives Eq.~\eqref{eq:graded-collapse}.

For the bit model, every Vandermonde entry has polynomial bit length in $L_j$;
Hadamard's inequality and Cramer's rule give polynomial bit length for its
inverse.  Because $g$ is fixed and the quotient basis of dimension
$\prod_jd_j$ is explicit, $\prod_j(L_j+1)$ is polynomial in the stated input
and $M$.  The polarization dimension $w^2$ is fixed.  Exact rational matrix
multiplication, tensor-product interpolation, quotient reduction, boundary
application, and division by $M!$ therefore use polynomially many bit
operations and produce polynomial-bit LC--AGP coefficients.  This conclusion
requires the rational embedding and quotient representation as input; it does
not claim an algorithm for discovering such an embedding from arbitrary
structure constants.

\bibliographystyle{apsrev4-2}
\bibliography{references/references}

\end{document}
